% 第9节：LSZ约化形式
\section{LSZ约化形式}
\label{sec:lsz}

\subsection{引言}

LSZ约化形式提供了一种系统的方法，用关联函数（Wightman函数）表示S-矩阵元。这对CTUFT中的微扰计算至关重要。

\subsection{约化公式推导}

\subsubsection{单粒子约化}

对于动量为$p$的入射粒子：
\begin{equation}
    |p, \text{in}\rangle = i \lim_{t \to -\infty} \int d^3x \, e^{-ip \cdot x} \overleftrightarrow{\partial_0} \hat{\mathcal{T}}(x) |0\rangle
\end{equation}

其中$f \overleftrightarrow{\partial_0} g = f \partial_0 g - (\partial_0 f) g$。

\subsubsection{S-矩阵元的约化}

\begin{theorem}[LSZ约化]
$2 \to 2$散射的连接S-矩阵元为：
\begin{multline}
    \langle p_3, p_4; \text{out} | p_1, p_2; \text{in} \rangle_C = \\
    i^4 \int \prod_{j=1}^4 d^4x_j \, e^{i(p_3 \cdot x_1 + p_4 \cdot x_2 - p_1 \cdot x_3 - p_2 \cdot x_4)} \\
    \times (\Box_1 + M_T^2)(\Box_2 + M_T^2)(\Box_3 + M_T^2)(\Box_4 + M_T^2) \\
    \times \langle 0 | T\{\mathcal{T}(x_1)\mathcal{T}(x_2)\mathcal{T}(x_3)\mathcal{T}(x_4)\} | 0 \rangle_C
\end{multline}
\end{theorem}

\subsection{推迟和超前函数}

约化公式涉及因果（Feynman）传播子。相关函数包括：

\begin{itemize}
    \item \textbf{推迟}：$G_R(x) = \theta(x^0) \langle [\mathcal{T}(x), \mathcal{T}(0)] \rangle$
    \item \textbf{超前}：$G_A(x) = -\theta(-x^0) \langle [\mathcal{T}(x), \mathcal{T}(0)] \rangle$
    \item \textbf{因果（Feynman）}：$G_F(x) = \langle T\{\mathcal{T}(x)\mathcal{T}(0)\} \rangle$
\end{itemize}

\subsection{谱表示}

\begin{theorem}[Källén-Lehmann谱表示]
两点函数有如下表示：
\begin{equation}
    G_F(p) = \int_0^\infty d\mu^2 \, \frac{\rho(\mu^2)}{p^2 - \mu^2 + i\epsilon}
\end{equation}
其中$\rho(\mu^2)$是谱密度，满足：
\begin{equation}
    \rho(\mu^2) = (2\pi)^3 \sum_n \delta^{(4)}(p - p_n) |\langle n | \mathcal{T}(0) | 0 \rangle|^2
\end{equation}
\end{theorem}

\subsection{在幺正性中的应用}

LSZ公式使得第\ref{sec:unitarity}节中的数值幺正性检验成为可能。光学定理：
\begin{equation}
    2 \, \text{Im} \mathcal{T}(p_1, p_2 \to p_1, p_2) = \sum_f |\mathcal{T}(p_1, p_2 \to f)|^2
\end{equation}
由关联函数的谱性质导出。
